\section{Explicit logarithmic evaluations and a rational dilogarithm value}
\label{jnew:sec:analytic}

The formal classification is complemented by full analytic evaluations.
These proofs use the positive real domain throughout, so their
logarithmic correction terms and multiples of $\pi^2$ are fixed exactly.

Define, as in Radchenko--Zagier~\cite{RadchenkoZagier},
\[
 \mathcal F(x)=F(x)-F(1)
  +\frac{\pi^2}{12}(x-2+x^{-1})-\frac14\log^2x,
 \qquad
 \mathcal J(x)=J(x)-\frac12\log^22
  +\frac{\pi^2}{24}(x-x^{-1}).
\]
Let $L(z)=\operatorname{Li}_2(z)+\tfrac12\log z\log(1-z)$ for
$0<z<1$, with the continuous endpoint values. The following are
established identities from Radchenko--Zagier,
equations (15) and (22a)~\cite{RadchenkoZagier}:
\begin{align}
 \mathcal F(x)-\mathcal F(x+1)-\mathcal F\left(\frac{x}{x+1}\right)
 &=L(1/2)-L\left(\frac{x}{x+1}\right),
 \label{jnew:eq:RZ-three}\\
 \mathcal F(x)+\mathcal F(1/x)&=0,
 \label{jnew:eq:RZ-recip}\\
 \mathcal F(2x)+\mathcal F(x/2)
 +\mathcal F\left(\frac{x+1}{2}\right)
 +\mathcal F\left(\frac{2x}{x+1}\right)-3\mathcal F(x)
 &=L\left(\frac{x}{x+1}\right)-L(1/2)
   +\tfrac12\log2\log x.
 \label{jnew:eq:RZ-Hecke2}
\end{align}
We use them as established analytic inputs, not as new results.
The Rogers five-term relation on $(0,1)^2$ is
\begin{equation}\label{jnew:eq:Rogers}
 L(u)+L(v)=L(uv)+L\left(\frac{u(1-v)}{1-uv}\right)
                    +L\left(\frac{v(1-u)}{1-uv}\right).
\end{equation}
Unlike its continuation to arbitrary real arguments, this displayed
identity holds as an equality in $\mathbb R$, with no unspecified
multiple of $\pi^2$.

The next lemma is a useful reformulation of known functional equations:
it also follows from Choie--Kumar's Theorem~3.2(1), equation (3.3),
after reciprocity and the three-term equation for
$\mathcal F$~\cite{ChoieKumar2023}. We give a direct positive-real
proof to fix every correction term. We do not claim it as a new
general functional equation. The explicit rational specializations
below are proposed additions to the manuscript; no exhaustive claim
of historical priority is made for those evaluations.

\begin{lemma}[A finite difference at adjacent positive arguments]
\label{jnew:lem:analytic-difference}
For every real $x>0$,
\begin{align}
 \mathcal J\left(\frac{x}{x+1}\right)
 ={}&\mathcal F(x)-\mathcal F(x+1)-\mathcal F(2x)
       +\mathcal F(2x+2)-L(1/2)\nonumber\\
 &\hspace{13mm}+\frac12\log2\log\left(\frac{x}{x+1}\right).
 \label{jnew:eq:analytic-difference}
\end{align}
\end{lemma}
\begin{proof}
Let $r=x/(x+1)$. By \eqref{jnew:eq:RZ-Hecke2} and
$\mathcal J(r)=\mathcal F(2r)-2\mathcal F(r)+\mathcal F(r/2)$,
\[
 \mathcal J(r)=\mathcal F(r)
 -\mathcal F\left(\frac{2x+1}{2x+2}\right)
 -\mathcal F\left(\frac{2x}{2x+1}\right)
 +L\left(\frac{x}{2x+1}\right)-L(1/2)
 +\tfrac12\log2\log r.
\]
Apply \eqref{jnew:eq:RZ-three} at $2x$ and $2x+1$. The two
middle $\mathcal F$ terms telescope to
\[
 -\mathcal F(2x)+\mathcal F(2x+2)+2L(1/2)
 -L\left(\frac{2x}{2x+1}\right)
 -L\left(\frac{2x+1}{2x+2}\right).
\]
In \eqref{jnew:eq:Rogers}, take
$u=2x/(2x+1)$ and $v=(2x+1)/(2x+2)$.
The three arguments on the right are $r$, $x/(2x+1)$, and $1/2$.
Thus all Rogers terms in the preceding expression reduce to $-L(r)$.
The result is
\[
 \mathcal J(r)=\mathcal F(r)-L(r)-\mathcal F(2x)
                       +\mathcal F(2x+2)+\tfrac12\log2\log r.
\]
One final application of \eqref{jnew:eq:RZ-three} at $x$ gives
\eqref{jnew:eq:analytic-difference}.
\end{proof}

\begin{theorem}[The complete consecutive rational family]
\label{jnew:thm:consecutive-analytic}
Put
\[
 S(n)=\sum_{k=1}^{n-1}\log^2\left(2\sin\frac{\pi k}{n}\right),
 \qquad S(1)=0.
\]
For every integer $n\ge1$,
\begin{align}
 \boxed{\displaystyle
 J\left(\frac{n}{n+1}\right)
 =\frac12\log^22
  -\frac{n^2-n-1}{24n(n+1)}\pi^2
  +\frac12\bigl(S(n)-S(n+1)-S(2n)+S(2n+2)\bigr).}
 \label{jnew:eq:consecutive-analytic}
\end{align}
Every logarithm has a positive algebraic argument.
\end{theorem}
\begin{proof}
The established integer value of $F$~\cite[eq.~(23)]{RadchenkoZagier} gives
\begin{equation}\label{jnew:eq:F-normalized-integer}
 \mathcal F(n)=\frac{n-1}{24}\pi^2
              +\frac14\log^2n+\frac12S(n).
\end{equation}
Insert this into \eqref{jnew:eq:analytic-difference}, with $x=n$.
The $\pi^2$ terms from the four $\mathcal F$ values sum to
$\pi^2/24$; subtracting $L(1/2)=\pi^2/12$ leaves $-\pi^2/24$.
The ordinary log-square terms sum to
$\tfrac12\log2\log((n+1)/n)$ and cancel the last term of
\eqref{jnew:eq:analytic-difference}. Therefore
\[
 \mathcal J\left(\frac n{n+1}\right)
 =-\frac{\pi^2}{24}
   +\frac12\bigl(S(n)-S(n+1)-S(2n)+S(2n+2)\bigr).
\]
Returning from $\mathcal J$ to $J$ proves the displayed formula.
\end{proof}

\begin{corollary}[Two small explicit values]\label{jnew:cor:small-evenodd}
With $\phi=(1+\sqrt5)/2$ and $\varepsilon=1+\sqrt2$,
\begin{align}
 J(3/4)&=-\frac{5\pi^2}{288}+\frac18\log^22
                         +\frac12\log^2\varepsilon,
 \label{jnew:eq:J34}\\
 J(4/5)&=-\frac{11\pi^2}{480}+\frac78\log^22
             +2\log^2\phi-\frac12\log^2\varepsilon.
 \label{jnew:eq:J45}
\end{align}
\end{corollary}
\begin{proof}
Direct pairing of the sine factors gives
\begin{align*}
 S(3)&=\tfrac12\log^23,& S(4)&=\tfrac32\log^22,\\
 S(5)&=\tfrac14\log^25+\log^2\phi,
 &S(6)&=\log^22+\tfrac12\log^23,\\
 S(8)&=\tfrac74\log^22+\log^2\varepsilon,
 &S(10)&=\log^22+\tfrac14\log^25+5\log^2\phi.
\end{align*}
For example, the logarithms of $2\sin(\pi/8)$ and
$2\sin(3\pi/8)$ have sum $\tfrac12\log2$ and difference
$\log\varepsilon$. Their squares therefore give the stated $S(8)$
formula. Substitute $n=3,4$ in
\eqref{jnew:eq:consecutive-analytic}.
\end{proof}

\begin{theorem}[A rational-dilogarithm value at $3/5$]
\label{jnew:thm:J35}
The following analytic identity holds:
\begin{align}
 \boxed{\displaystyle
 J(3/5)=\operatorname{Li}_2(1/3)-\frac12\operatorname{Li}_2(1/5)
 -\frac{7\pi^2}{180}+\frac12\log^22+\frac12\log^23
 -\frac14\log^25+\log^2\phi.}
 \label{jnew:eq:J35}
\end{align}
Its rational-dilogarithm core is exactly
$[1/3]-\tfrac12[1/5]$ in the formal calculus.
\end{theorem}
\begin{proof}
Take $x=3/2$ in \eqref{jnew:eq:analytic-difference} and use
\eqref{jnew:eq:RZ-recip}:
\[
 \mathcal J(3/5)
 =-\mathcal F(2/3)+\mathcal F(2/5)-\mathcal F(3)+\mathcal F(5)
  -L(1/2)+\tfrac12\log2\log(3/5).
\]
Equation \eqref{jnew:eq:RZ-three} at $x=2$ says
$\mathcal F(2/3)=\mathcal F(2)-\mathcal F(3)-L(1/2)+L(2/3)$.
Consequently,
\begin{equation}\label{jnew:eq:J35-short-proof}
 \mathcal J(3/5)=\mathcal F(2/5)+\mathcal F(5)-\mathcal F(2)
       -L(2/3)+\tfrac12\log2\log(3/5).
\end{equation}
Use the established Radchenko--Zagier evaluation
\cite[eq.~(24)]{RadchenkoZagier}
\begin{align*}
 F(2/5)-F(1)
  ={}&\tfrac12\operatorname{Li}_2(4/5)-\frac{\pi^2}{5}
    +\log^2\left(2\sin\frac\pi5\right)
    +\log^2\left(2\sin\frac{2\pi}5\right)
    +\log2\log(2/5),
\end{align*}
together with \eqref{jnew:eq:F-normalized-integer} at $2,5$.
The two log-sine squares in this formula have sum
$\tfrac18\log^25+\tfrac12\log^2\phi$.
Simplifying \eqref{jnew:eq:J35-short-proof} and applying Euler's
reflection identity to $\operatorname{Li}_2(4/5)$ and
$\operatorname{Li}_2(2/3)$ gives
\[
 \mathcal J(3/5)=\operatorname{Li}_2(1/3)
    -\tfrac12\operatorname{Li}_2(1/5)-\frac{\pi^2}{12}
    +\tfrac12\log^23-\tfrac14\log^25+\log^2\phi.
\]
Since $3/5-5/3=-16/15$, the conversion from $\mathcal J$ to $J$
adds $\tfrac12\log^22+2\pi^2/45$, proving
\eqref{jnew:eq:J35}. All Rogers and Euler reflection arguments
lie in $(0,1)$, so there is no branch ambiguity.

For an independent boundary check,
\[
 \partial\left([1/3]-\tfrac12[1/5]\right)
 =\langle2\rangle\wedge\langle3/5\rangle,
\]
which agrees with \eqref{jnew:eq:odd-boundary} because the symbols
at conductors $3$ and $5$ vanish.
\end{proof}

\begin{proposition}[Reciprocity and the remaining odd rational cores]
\label{jnew:prop:reciprocity}
For every $x>0$,
\begin{equation}\label{jnew:eq:analytic-reciprocity}
 J(x)+J(1/x)=\log^22.
\end{equation}
In particular, the established $J(1/3)$ and $J(1/5)$ formulas yield
\begin{align*}
 J(3)&=-\frac{\pi^2}{9}+\frac12\log^22+\frac12\log^23
                                     +\operatorname{Li}_2(1/3),\\
 J(5)&=-\frac{7\pi^2}{60}+\frac12\log^22+\frac14\log^25
                    +\log^2\phi+\frac12\operatorname{Li}_2(1/5).
\end{align*}
The value $J(5/3)$ is $\log^22-J(3/5)$. These account for all
nontrivial odd/odd cases in Theorem~\ref{jnew:thm:classification}.
\end{proposition}
\begin{proof}
Integration by parts gives
\[
 J(x)=\log^22-
 x\int_0^1\frac{t^{x-1}\log(1+t)}{1+t^x}\,dt.
\]
Substitute $u=t^x$ in the remaining integral to obtain $J(1/x)$.
The boundary at zero vanishes and the boundary at one is $\log^22$.
This reciprocity is elementary and already known; see, for example,
Choie--Kumar~\cite[Theorem~3.1]{ChoieKumar2023}. It is included to
make the analytic consequences complete. The displayed values follow
by substitution in the existing $1/3$ and $1/5$ formulas.
\end{proof}
